% !TeX root = ../../Book5.tex
% 纲要标题：### 21.5 型、对合与古典群
% 纲要位置：工作记录/计划纲要.md，第 4415 行。
% 本轮补充的纲要细目：
% ---
% 写作范围：具体型、对合及四元数模型；不外推为一般扭曲形式分类。

\section{型、对合与古典群}
\label{b5:sec:2105}

% 纲要细目：
% * 从非退化型的分裂矩阵模型解释正交、辛、酉群及相似变换群，写清底域、特征和型的约定
% * 区别等距群、相似变换群、代数自同构群与中心商；代数群满态射不保证有理点映射满射
% * 选一个四元数的非分裂模型，并重述所需代数事实；不以一般扭曲形式分类或未写的非交换上同调作为前提
% #### 21.5.1 等距群、相似变换群与乘子
% #### 21.5.2 带对合代数的自同构群
% #### 21.5.3 分裂模型与四元数例子

正交群、辛群和酉群最初由保持型的矩阵定义。把型换成端同态代数上的伴随对合以后，同一批方程也能在非分裂代数中使用。本节先固定具体模型，再比较等距、相似变换和代数自同构；它们之间有联系，却不是同一个群。

\subsection{等距群、相似变换群与乘子}
\label{b5:subsec:2105:01}

设 \(F\) 为特征不为二的域，取 \(K=F\) 或一个可分二次扩张 \(K/F\)，并以 \(a\mapsto\bar a\) 分别表示恒等或非平凡共轭。给定 \(H\in\operatorname{GL}_n(K)\)、\(n\ge1\)，假设
\(H^*=\epsilon H\)、\(\epsilon\in\{1,-1\}\)，其中 \(X^*=\bar X^{\mathsf T}\)。型为
\(b(u,v)=\bar u^{\mathsf T}Hv\)，伴随对合为
\(\sigma_H(X)=H^{-1}X^*H\)。
\begin{definition}\label{b5:def:classical-similitude}
设 \(F\) 为特征不为二的域，\(K=F\) 或 \(K/F\) 为可分二次扩张，共轭分别为恒等或非平凡域自同构。取 \(H\in\operatorname{GL}_n(K)\)、\(n\ge1\)，满足 \(H^*=\epsilon H\)、\(\epsilon\in\{1,-1\}\)，其中 \(X^*=\bar X^{\mathsf T}\)。对非退化型 \(b(u,v)=\bar u^{\mathsf T}Hv\)，定义等距群与相似变换群
\[
\begin{aligned}
\operatorname{Isom}(b)&=\Bigl\{g\in\operatorname{GL}_n(K)\mid g^*Hg=H\Bigr\},\\
\operatorname{Sim}(b)&=\Bigl\{g\in\operatorname{GL}_n(K)\mid
g^*Hg=cH\text{ 对某个 }c\in F^\times\Bigr\}.
\end{aligned}
\]
相似变换的这个唯一标量 \(c\) 称为乘子，记为 \(\mu(g)\)。当 \(K=F\)、\(H\) 对称时，两群分别记为 \(\operatorname O(b)\)、\(\operatorname{GO}(b)\)；当 \(H\) 交替时，记为 \(\operatorname{Sp}(b)\)、\(\operatorname{GSp}(b)\)；二次扩张上的 Hermitian 情形记为 \(\operatorname U(b)\)、\(\operatorname{GU}(b)\)。
\end{definition}
这些式子不仅定义 \(F\) 上的抽象群。对任意交换 \(F\) 代数 \(R\)，把系数扩大到 \(K\otimes_FR\)，令共轭固定 \(R\)，便得到相同的矩阵方程；将非零乘子换为 \(R^\times\) 中的元，就定义相应仿射群概形。等距群是一般线性群的闭子群，相似群可在 \(\operatorname{GL}_n(K)\times\mathbb G_m\) 中用方程实现；\(K/F\) 有限维使这些都是有限个 \(F\) 坐标方程。
\begin{proposition}\label{b5:prop:classical-multiplier}
设 \(F\) 为特征不为二的域，\(K=F\) 或 \(K/F\) 为可分二次扩张，共轭分别为恒等或非平凡域自同构，\(X^*=\bar X^{\mathsf T}\)。取 \(n\ge1\)、\(H\in\operatorname{GL}_n(K)\)，满足 \(H^*=\epsilon H\)、\(\epsilon=\pm1\)。令 \(b\) 为 \(K^n\) 上的非退化型 \(b(u,v)=\bar u^{\mathsf T}Hv\)，\(\sigma_H(X)=H^{-1}X^*H\)，记 \(\operatorname{Sim}(b)\)、\(\operatorname{Isom}(b)\) 分别为相似变换群和等距群。则乘子给出群同态 \(\mu:\operatorname{Sim}(b)\to F^\times\)，其核为 \(\operatorname{Isom}(b)\)。并且
\[
g\in\operatorname{Sim}(b),\ \mu(g)=c
\quad\Longleftrightarrow\quad \sigma_H(g)g=cI.
\]
\end{proposition}
\begin{proof}
非退化型至少有一个非零配对值，故乘子唯一。两个相似变换连续作用使配对值乘两个标量的积，逆变换乘以逆标量，所以是同态，乘子一恰是等距条件。最后左乘 \(H^{-1}\) 即将矩阵方程改写为所示对合方程。
\end{proof}
乘子的实际像须另外判断。在正定实型上它只能为正数，且 \(\sqrt c I\) 实现每个正数；在分裂辛型
\(J=\begin{psmallmatrix}0&I_m\\-I_m&0\end{psmallmatrix}\)
上，\(\operatorname{diag}(cI_m,I_m)\) 实现每个 \(c\in F^\times\)。所以写核与像时不能一概补上“到 \(F^\times\) 满射”。

\subsection{带对合代数的自同构群}
\label{b5:subsec:2105:02}

\begin{proposition}\label{b5:prop:involution-automorphism-similitude}
设 \(F\) 为特征不为二的域，\(K=F\) 或 \(K/F\) 为可分二次扩张，共轭分别为恒等或非平凡域自同构，\(X^*=\bar X^{\mathsf T}\)。取 \(H\in\operatorname{GL}_n(K)\)、\(n\ge1\)，满足 \(H^*=\epsilon H\)、\(\epsilon=\pm1\)。令 \(A=M_n(K)\)、\(b(u,v)=\bar u^{\mathsf T}Hv\) 为 \(K^n\) 上的非退化型、\(\sigma(X)=H^{-1}X^*H\)，\(\operatorname{Sim}(b)\) 为型 \(b\) 的相似变换群。保持 \(\sigma\) 的 \(K\) 代数自同构组成的群记为
\(\operatorname{Aut}_K(A,\sigma)\)。内自同构
\(\operatorname{Int}(g)(X)=gXg^{-1}\)
给出群同构
\[
\operatorname{Sim}(b)/K^\times
\ \cong\ \operatorname{Aut}_K(A,\sigma),
\]
其中 \(K^\times\) 作为非零标量矩阵嵌入相似群，其乘子为 \(\bar a a\)。
\end{proposition}
\begin{proof}
矩阵代数的每个 \(K\) 代数自同构都是内自同构，这是第二卷定理17.3.1的 Skolem--Noether 结论在中心单代数 \(M_n(K)\) 上的应用。直接计算
\[
\sigma(gXg^{-1})=\sigma(g)^{-1}\sigma(X)\sigma(g).
\]
它等于 \(g\sigma(X)g^{-1}\) 对所有 \(X\) 成立，当且仅当
\(\sigma(g)g\) 与所有矩阵交换，即为某个 \(cI\)、\(c\in K^\times\)。对该等式再用 \(\sigma\)，得到 \(\bar c=c\)，故 \(c\in F^\times\)。由前一命题，这正是相似条件。内自同构的核是矩阵代数的中心单位 \(K^\times\)，而这些标量全属于相似群，遂得所示同构。
\end{proof}
这里比较的是指定域上的点群，并且自同构要求逐点固定 \(K\)。二次扩张上若允许自同构在中心上也作共轭，会得到另一个群。对一般带对合中心单代数，相同的计算仍把内自同构与条件 \(\sigma(g)g\in F^\times\) 联系起来；非分裂情形中的型及对合细节见第二卷附录B。
\ProseParagraphBreak
还须区分代数群满态射与域上点的满射。复数上，标量平方根能把任意可逆二阶矩阵缩放为行列式一矩阵，所以
\(\operatorname{SL}_2(\mathbb C)\to\operatorname{PGL}_2(\mathbb C)\)
满射。相应代数群态射在实数上也定义良好，但实点的像只含行列式为正的射影矩阵类：缩放把行列式乘一个正平方，不可能把负行列式改成一。\(\operatorname{diag}(-1,1)\) 的射影类便没有实的 \(\operatorname{SL}_2\) 原像。

\subsection{分裂模型与四元数例子}
\label{b5:subsec:2105:03}

在 \(F=\mathbb R\) 上，Hamilton 四元数代数
\[
\mathbb H=\mathbb R\oplus\mathbb Ri\oplus\mathbb Rj\oplus\mathbb Rk,
\quad i^2=j^2=k^2=-1,\quad ij=k=-ji
\]
带有共轭 \(\bar q\) 和范数
\(N(q)=q\bar q=q_0^2+q_1^2+q_2^2+q_3^2\)。
共轭反转乘法，故 \(N(uv)=N(u)N(v)\)。于是
\[
\operatorname{SL}_1(\mathbb H)(\mathbb R)=\Bigl\{q\in\mathbb H\mid N(q)=1\Bigr\}
\]
由一个四元二次方程定义为实代数群；逆由共轭给出，不需要除以变量。
\ProseParagraphBreak
复化后它变成分裂模型。把 \(i,j\) 分别送到
\[
\begin{pmatrix}\mathrm i&0\\0&-\mathrm i\end{pmatrix},
\qquad
\begin{pmatrix}0&1\\-1&0\end{pmatrix}.
\]
两矩阵平方为 \(-I\)、反交换，并且它们连同单位及乘积线性无关，所以给出
\(\mathbb H\otimes_{\mathbb R}\mathbb C\cong M_2(\mathbb C)\)。
计算一般元素的行列式恰为 \(N(q)\)，因此复化后的群为 \(\operatorname{SL}_2\)。在实点上，四元数单位球面是有界闭集，而 \(\operatorname{SL}_2(\mathbb R)\) 含任意大的 \(\operatorname{diag}(t,t^{-1})\)；这两个实模型的差异不由其共同的复化消除。
\begin{example}\label{b5:exm:quaternion-conjugation-group}
单位四元数通过 \(x\mapsto qxq^{-1}\) 作用于纯虚三维空间
\(\operatorname{Im}\mathbb H\)。它保持该空间，因为实部是 \((x+\bar x)/2\)，且共轭作用保持实部；它也保持正定范数。作用的核是 \(\{\pm1\}\)：若与 \(i,j,k\) 都交换，便在 \(\mathbb H\) 的中心 \(\mathbb R\) 中，再加范数一就只剩两个元。
例如 \(q=i\) 固定 \(i\) 并把 \(j,k\) 变号；而
\(q=(1+i)/\sqrt2\) 固定 \(i\)，将 \(j\) 送到 \(k\)、\(k\) 送到 \(-j\)。这是由非交换代数乘法产生的具体正交作用。
\end{example}
分裂矩阵型、非分裂四元数、等距群与代数自同构因而可以互相解释。它们也提醒我们：复数域上的根系或李代数类型只记录部分结构；改换基域、对合或中心商以后，点群与表示下降问题还需要重新核对。
